\documentclass[12pt,a4paper]{article}

\usepackage[T1]{fontenc}
\usepackage[utf8]{inputenc}
\usepackage[polish]{babel}
\usepackage{amsmath,amssymb}
\usepackage{enumitem}
\usepackage[a4paper,margin=2.5cm]{geometry}


\begin{document}
\begin{center} \bf Tematy prac magisterskich 2026/2027 \end{center}
\begin{itemize}
\item
  Asymptotyczna zbieżność i geometryczna ergodyczność
  łańcuchów Markowa w bayesowskiej analizie danych

\item
  Modelowanie xiągłem zmienności stochastycznej przy użyciu procesu
COGARCH.

\item
  Wybrane metody estymacji krzywej uczenia.

\item
  Nieskończona wersja gry w kółko i krzyżyk. 

\item
  {\bf Zastosowania twierdzenia Ramseya}

Teoria Ramseya zajmuje się zjawiskiem pojawiania się uporządkowanych struktur
w dostatecznie dużych układach, niezależnie od sposobu ich podziału lub
kolorowania. Klasycznym przykładem są \emph{liczby Ramseya} $R(k,l)$,
określające najmniejszą liczbę $n$, dla której każde dwukolorowanie krawędzi
grafu pełnego $K_n$ zawiera albo czerwony podgraf $K_k$, albo niebieski
podgraf $K_l$.

Celem pracy będzie przedstawienie wybranych zastosowań liczb Ramseya i
związanych z nimi metod. Szczególna uwaga może zostać poświęcona
interpretacjom grafowym i kombinatorycznym oraz zastosowaniom teorii Ramseya
do problemów, w których chodzi o wykazanie, że w dostatecznie dużym zbiorze
obiektów musi wystąpić określona regularność. Możliwe jest również omówienie
zastosowań w teorii grafów, kombinatoryce, teorii liczb oraz informatyce
teoretycznej.

\section*{Przykładowe zagadnienia}

\begin{itemize}[leftmargin=2em]
    \item liczby Ramseya $R(k,l)$ i ich podstawowe własności;
    \item zastosowania teorii Ramseya w teorii grafów;
    \item problemy ekstremalne i gwarantowanie występowania struktur
          monochromatycznych;
    \item zastosowania w kombinatoryce i teorii liczb;
    \item wybrane wyniki dotyczące liczb Ramseya dla grafów i innych
          struktur kombinatorycznych.
\end{itemize}

\section*{Literatura}

\begin{enumerate}[leftmargin=2em]
\item K.\ Hrbacek, T.\ Jech,
    \emph{Introduction to Set Theory}, Marcel Dekker, New York, 1999.

\item R.\ L.\ Graham, B.\ L.\ Rothschild, J.\ H.\ Spencer,
    \emph{Ramsey Theory}, 2nd ed., John Wiley \& Sons, New York, 1990.

\item Wikipedia contributors,
    \emph{Twierdzenie Ramseya}, 2026.

\item C.\ Franks,
    \emph{Ramsey Theory, Probabilistic Method: Math 454 Lecture 16}, 2017.

\item R.\ Żelazko,
    \emph{Charakterystyka rozbicia zbioru co najwyżej przeliczalnego}, 2013.

\item H.\ M.\ Wong, M.\ H.\ Leung, W.\ Y.\ Wong, H.\ K.\ Hui, T.\ C.\ Mak,
    \emph{The Erd{\H{o}}s--Szekeres Conjecture (``Happy End Problem'')},
    Hang Lung Mathematics Awards 4, 271--298, 2010.

  \item R.\ L.\ Graham, B.\ L.\ Rothschild, J.\ H.\ Spencer,
    \emph{Ramsey Theory}, 2nd ed., Wiley-Interscience, 1990.

    \item R.\ Diestel,
    \emph{Graph Theory}, 5th ed., Springer, 2017.
    --- rozdziały dotyczące teorii Ramseya i zastosowań w teorii grafów.

    \item P.\ Erd\H{o}s, G.\ Szekeres,
    A combinatorial problem in geometry,
    \emph{Compositio Mathematica} \textbf{2} (1935), 463--470.

    \item J.\ H.\ van Lint, R.\ M.\ Wilson,
    \emph{A Course in Combinatorics}, 2nd ed., Cambridge University Press,
    2001.

    \item S.\ Jukna,
    \emph{Extremal Combinatorics: With Applications in Computer Science},
    2nd ed., Springer, 2011.

    \item G.\ Chartrand, P.\ Zhang,
    \emph{A First Course in Graph Theory}, Dover Publications, 2012.
  \end{enumerate}
  \end{itemize}

\end{document}
